\section{Introduction}\label{sec:intro}

Let \(q:X\to Q\) be a quotient map, let \(F:X\to Y\) be a map one
wants to transport through that quotient, and let \(r:Y\to R\) be the
target identification under which the transported output is to be
read.  The elementary question is whether the composite \(rF\)
depends only on the \(q\)-class of the input.  Equivalently, does
there exist a map \(\overline F:Q\to R\) with
\[
  \overline F q = rF?
\]
If \(x\) and \(x'\) have the same image in \(Q\) but \(rF(x)\ne
rF(x')\), then no such \(\overline F\) can exist.  Conversely, when
\(q\) is surjective, the absence of such pairs is exactly the
condition for descent.  There are two canonical repairs.
One may refine the source quotient so that the offending pair is no
longer identified, or one may coarsen the target by quotienting \(R\)
by the smallest equivalence relation generated by the identifications
\[
  rF(x)\sim rF(x')
  \qquad\text{whenever}\qquad
  q(x)=q(x').
\]
The latter repair forgets precisely the distinction that the source
quotient cannot support.  The obstruction is thus not merely a
failure: it is a diagnostic object, and the minimal repairs are
determined by the same data that produced it.

\paragraph{A running example.} Take \(X=\{1,2,3,4\}\) and
\(Q=\{a,b,c\}\) with \(q(1)=q(2)=a\), \(q(3)=b\), \(q(4)=c\); take
\(Y=R=\{u,v,w\}\) with \(r\) the identity, and \(F(1)=u\),
\(F(2)=v\), \(F(3)=F(4)=w\).  The pair \((1,2)\) is split:
\(q(1)=q(2)\) but \(F(1)\ne F(2)\), so \(F\) does not descend through
\(q\).  The source repair separates \(1\) and \(2\); the target repair
identifies \(u\) and \(v\).  Each chapter opening below returns to this
example briefly.

The example uses no operator theory, no regularity hypothesis and no
category larger than the quotient construction itself.  Yet the same
shape appears in more elaborate settings.  A local family may agree on
overlaps but fail to come from a global object.  A symmetric
construction may be unable to choose a branch without adding
symmetry-breaking data.  A formally closed record may exist without
being physically realised.  A finite observation may determine all
declared future probes, or fail to, precisely because some future probe
separates states that the present quotient cannot see.  In each case
there is a setup, an obstruction, a repair or boundary, and a scope set
by the hypotheses.

\paragraph{What the paper does.}
\FIV\ is a catalog of such patterns.  It introduces no new foundation;
it states, as theorems with explicit hypotheses, \NumLaws\ structural
laws that recur across the Six Birds series, and proves them.  Each law
is presented in the same order: its setup, the theorem, the proof, a
case enumeration when the statement branches, and, where useful, a
paragraph saying what the law does not claim
(\Cref{def:normal-form}).  The uniform presentation lets a reader ask
the same questions of every entry: what are the objects, what is the
obstruction, what is the repair or classification, and what lies
outside the claim?  It is not meant to suggest that the laws are one
theorem.

Most laws are elementary.  Many are the descent criterion applied to
particular data, and several carry much of their content in their
hypotheses: the domination bounds of Duality Fixity (\Fref{F14}) exist
exactly when the readout they confine vanishes almost everywhere, and part (i) of the
Hiddenness Normal Form (\Fref{F13a}) assumes four of its five
conclusions.  The value of the catalog lies in stating each pattern
once, with its hypotheses and its scope visible, so that it can be
recognized and checked in a new setting.

A \emph{law} here is not a physical law or a principle detached from
hypotheses.  It is a named theorem stated at the level of the data in
its setup.  Some laws are quotient-factorization statements, some are
finite classifications, and some are boundary statements showing that a
desirable implication fails without stronger data.  The names of the
laws, such as Probability, Causality or Entropy, say what kind of data
the maps carry; they do not announce new theorems about probability,
causality or entropy.

Most laws apply one criterion to different data: a readout descends
through a quotient exactly when its split-pair obstruction, the set of
pairs that the quotient merges and the readout separates, is empty
(\Cref{def:split-pair}).  A few use arguments of other kinds: a
Loewner-order inequality (No-Needles Stability, \Fref{F6}), a trace
squeeze (Duality Fixity, \Fref{F14}), fiber weights (\Fref{F23}),
generating relations (\Fref{F27}, \Fref{F29}), a well-founded minimum
(\Fref{F42}), connectedness of \([0,1]\) (\Fref{F44}), and finite
counterexamples in the closure--reality no-go theorems.  Several laws,
such as Local--Global Obstruction (\Fref{F4}), also give a
classification into mutually exclusive cases.

\paragraph{Inherited apparatus.}
The paper uses three earlier papers of the series as background, in
the compact form recorded in \Cref{sec:apparatus}.  Foundations~I
supplies the language of closure formation, the completion of a
structure under declared rules \citep{Tsiokos2026EmergenceCalculus};
Foundations~II supplies admissibility and quotient packaging, the
standard act of retaining only the equivalence classes of a structure
\citep{Tsiokos2026FoundationsII}; and Foundations~III supplies a finite
audited interaction judgment, which records when a proposed interaction
is legitimately visible to the apparatus
\citep{Tsiokos2026FoundationsIII}.  Two rules of those papers enter
proofs here as explicit premises: the no-overread rule (NO) in
\Fref{F11} and \Fref{F12}, and the residual-coverage rule (RC) in
\Fref{F9}.

\paragraph{Organization.}
The laws are cited by F-codes, which are fixed identifiers and are not
in document order; the map of results in \Cref{subsec:reading-guide}
gives the printed number and page of each.  Six chapters group the laws
by the kind of question they answer.  The Access chapter
(\Cref{sec:access}) studies what an instrument can distinguish; its
first theorem says that an observation signature determines a canonical
access quotient.  The Stability chapter (\Cref{sec:stability}) studies
when declared future behavior is already controlled by present data and
when residuals must be accounted for.  The Transport chapter
(\Cref{sec:transport}) is the easiest entry point, since its laws are
closest to standard descent, gluing and quotient comparison.  The
Symmetry chapter (\Cref{sec:symmetry}) treats obstructions coming from
involutions, branch selection and anomalies.  The Status, Records and
Coherence chapter (\Cref{sec:status-records-coherence}) collects laws
in which records, facts, persistence, explanation and counterfactuals
are part of the mathematical object.  The Self-Reference and Limits
chapter (\Cref{sec:self-reference-limits}) records horizon,
information, singular-limit, locality, topology and completeness
boundaries.  A final chapter (\Cref{sec:anchor:closure-reality-boundary})
separates formal closure from bare physical realisability.
\Cref{sec:substrate-instances} maps the laws to the papers of the series
that instantiate them, \Cref{sec:sketches} lists candidates that are not
yet laws, and \Cref{sec:nonclaims} collects the scope limits of the
catalog.

\paragraph{Status of the results.}
Every law is proved here, with two qualifications.  The Hiddenness
Normal Form (\Fref{F13a}) is a conditional theorem: it assumes, as an
explicit hypothesis, surjectivity of the current quotient and four of
its five conclusions, and proves the fifth together with the
equivalence of predictive surplus with failure of predictive collapse.
The hypothesis is motivated by the saturated-SAT construction of
\citet{Tsiokos2026PvNP} but is not verified for it.  The
Closure--Reality Boundary resolution (\Fref{F53}) imports one
source-accounting equivalence from the paper on audited operational
realisability \citep{Tsiokos2026AOR}; its two no-go theorems are proved
here.  The substrate papers cited here are papers of the series by the same
author.  We say that a law is \emph{anchored} to a substrate paper when
it generalizes a theorem of that paper (\Fref{F6}, \Fref{F14}) or
imports one, explicitly identified (clause (b) of \Fref{F53}), and \emph{motivated} by one when that paper suggests its
hypotheses without establishing them (\Fref{F13a}, \Fref{F13b});
statements about what the substrate papers prove are those papers' own
claims.  \Cref{tab:substrate-instances} lists these connections.  Applications to particular
mathematical, physical or engineered systems appear after each law as
\emph{layer instantiations}, each marked as a special case or as an
analogy (\Cref{subsec:reading-guide}); no proof depends on them.  The
catalog is not exhaustive: \Cref{sec:sketches} records candidates that
remain sketches.  The results are formalized in Lean~4;
\Cref{app:formalization} describes the correspondence.

\paragraph{Reading path.}
We suggest reading this introduction and \Cref{sec:apparatus}, then
\Cref{sec:transport}, which depends only on \Cref{sec:apparatus} and
can be read before \Cref{sec:access,sec:stability}; the remaining
chapters can be read in any order.  Readers familiar with the earlier
foundations papers may skim \Cref{sec:apparatus} for notation and use
the summary table at the end of each chapter.  Readers coming from a
substrate paper should start from the corresponding row of
\Cref{tab:substrate-instances}.
